\subsection{Preesquemas sobre una base regular de dimensión 1; clausura de un
subpreesquema cerrado de la fibra genérica.}
\label{subsection:IV.2.8-fr}

\begin{proposition}[2.8.1]
\label{IV.2.8.1-fr}
Sean $Y$ un preesquema localmente noetheriano, regular, irreducible y de
dimensión 1, cuyo punto genérico es $\eta$, $f:X\to Y$ un morfismo,
$X_\eta=f^{-1}(\eta)$ la fibra en el punto

\oldpage[IV-2]{34}
genérico, $i:X_\eta\to X$ el morfismo canónico. Sean $\mathcal F$ un
$\mathcal O_X$-módulo cuasicoherente, $\mathcal F_\eta=i^*(\mathcal F)$,
$\mathcal G'$ un $\mathcal O_{X_\eta}$-módulo cociente de $\mathcal F_\eta$,
y sea $\overline{\mathcal G'}$ el $\mathcal O_X$-módulo imagen de $\mathcal F$
por el homomorfismo compuesto $(\mathbf{0}_{\mathrm I},4.4.3.2)$
\[
\mathcal F\xrightarrow{\rho_{\mathcal F}}i_*(i^*(\mathcal F))
\longrightarrow i_*(\mathcal G').
\]
Entonces $\overline{\mathcal G'}$ es un $\mathcal O_X$-módulo cociente de
$\mathcal F$, cuasicoherente y plano respecto de $f$, tal que
$i^*(\overline{\mathcal G'})=\mathcal G'$, y es el único
$\mathcal O_X$-módulo cociente de $\mathcal F$ que tiene estas propiedades.
\end{proposition}

\begin{proof}
Como $i$ es cuasicompacto y cuasiseparado (1.1.2 y 1.2.2), se deduce de
(1.7.4) que, para todo $\mathcal O_{X_\eta}$-módulo cuasicoherente
$\mathcal H'$, $i_*(\mathcal H')$ es un $\mathcal O_X$-módulo
cuasicoherente; además, para todo abierto $U$ de $X$, se tiene
$(i_*(\mathcal H'))|U=(i|U\cap X_\eta)_*(\mathcal H'|U\cap X_\eta)$
por definición $(\mathbf{0}_{\mathrm I},3.4.1)$. Si se demuestra la proposición cuando
$X$ e $Y$ son afines, se deducirá por pegado en el caso general, vista la
afirmación de unicidad válida en el caso afín. En otras palabras, la cuestión se reduce
a demostrar el

\begin{lemma}[2.8.1.1]
\label{IV.2.8.1.1-fr}
Sean $A$ un anillo noetheriano regular $(\mathbf{0},17.3.6)$, íntegro de dimensión
1, $K$ su cuerpo de fracciones, $M$ un $A$-módulo, $N'$ un $K$-módulo
cociente de $M_{(K)}=M\otimes_A K$ por un sub-$K$-módulo $P'$,
$\overline{N'}$ la imagen de $M$ por el homomorfismo compuesto
$M\to M_{(K)}\to N'$. Entonces $\overline{N'}$ es un $A$-módulo plano y es
el único módulo cociente $N$ de $M$ que sea un $A$-módulo plano y tal que el
núcleo del homomorfismo sobreyectivo $M_{(K)}\to N_{(K)}$ sea igual a $P'$.
\end{lemma}

Como, para todo ideal maximal $\mathfrak m$ de $A$, $A_{\mathfrak m}$ es un
anillo local regular de dimensión 1 y, por tanto, un anillo de valoración
discreta, es lo mismo decir que un $A$-módulo $N$ es plano que decir que es
\emph{sin torsión} $(\mathbf{0}_{\mathrm I},6.3.4)$. Como $N'$ es un espacio vectorial
sobre $K$, es un $A$-módulo sin torsión, por lo que lo mismo ocurre con
$\overline{N'}$, submódulo de $N'$; además, es inmediato comprobar que
$\overline{N'}_{(K)}$ se identifica con $N'$. Recíprocamente, si $N$ es un
$A$-módulo cociente de $M$ que tiene las propiedades del enunciado, el hecho de
que $N$ sea un $A$-módulo plano implica que el homomorfismo canónico
$N\to N_{(K)}=N\otimes_A K$ es \emph{inyectivo}. Como $N_{(K)}$ se identifica con
$N'$, la conclusión se deduce de la conmutatividad del diagrama
\[
\xymatrix{
M \ar[r] \ar[d] & N \ar[d] \\
M_{(K)} \ar[r] & N_{(K)}
}
\]
\end{proof}
